%% =========================================================
%% Point-by-point response to reviewers -- CHB-D-26-04350
%% Standalone; compile from the repo root AFTER manuscript.tex has been
%% compiled (this file imports manuscript.aux via \externaldocument so that
%% \ref{tab:...}/\ref{sec:...} resolve to the manuscript's own numbers).
%%   pdflatex response_to_reviewers && pdflatex response_to_reviewers
%%
%% Left-hand "Reviewer comment" cells quote the reviewer's verbatim wording
%% (elisions marked with [...]); Q#/P# labels follow the review letter.
%% All numbers synced to the full-data answers-only tables (2026-07-17);
%% cross-check against analysis/output_tables/ before submission.
%% =========================================================
\documentclass[11pt]{article}
\usepackage[margin=1in]{geometry}
\usepackage{array}              % \newcolumntype, \arraybackslash
\usepackage{booktabs}
\usepackage{longtable}          % tables that break across pages
\usepackage{amsmath}
% hyperref intentionally omitted: it is only used for cosmetic (hidden) ref/cite
% links here, cross-PDF links to the separate manuscript do not work anyway, and
% plain xr + hyperref breaks external \ref inside the longtable cells
% ("Paragraph ended before \Hy@setref@link was complete"); xr-hyperref is not
% installed in this TeX Live. Plain \ref/\cite render the external numbers fine.
\usepackage{xr}                 % cross-reference the manuscript's labels
\externaldocument{manuscript}
\setlength{\parskip}{6pt}
\setlength{\parindent}{0pt}
\renewcommand{\arraystretch}{1.15}

\newcolumntype{L}{>{\raggedright\arraybackslash}p{0.30\textwidth}}
\newcolumntype{R}{>{\raggedright\arraybackslash}p{0.62\textwidth}}

\begin{document}

\begin{center}
{\large\bfseries Response to Reviewers}\\[2pt]
Manuscript CHB-D-26-04350\\
\emph{Help Converts Newcomers, Not Veterans: Generalized Reciprocity and Platform Engagement on Stack Overflow}
\end{center}

We thank the editor and both reviewers for their careful and constructive reading of the manuscript. In preparing this revision we substantially re-analysed the data: we re-matched treated and control questions using nearest-neighbor propensity-score matching \emph{with replacement} (caliper 0.05), with exact strata on calendar year, primary tag, and the number of tags, and re-fit every Cox model on the resulting sample, so that all reported estimates derive from a single, consistent pipeline run. The principal changes are: (i) matched-pair bootstrap confidence intervals for the pooled treatment effect, replacing the earlier claim that clustering was unnecessary; (ii) a full decomposition of the helping outcome by help type (answers, comments, edits, and their composites), which retains answers --- now explicitly justified on construct grounds --- as the headline outcome and shows the reciprocity signal is answer-borne; (iii) a flexible, non-parametric response-time specification (discrete bins), re-estimated with observable answer-quality controls; (iv) new robustness analyses for observable selection, entry-cohort heterogeneity, and left-truncation/churn; (v) a more modest, tenure-focused framing throughout, under which the pooled association is small, the newcomer gradient is the most robust pattern, the response-time result is offered as qualified support for an immediacy account, and the quality-controlled estimates are read as a conservative post-treatment bound; (vi) a shortened Theory section that is linked directly to the empirical design, with interpretation-relevant appendix material promoted into the main text; and (vii) a full language and consistency edit of the revised manuscript. One further transparency item: a legacy supplementary analysis of lifetime contributions was removed from the appendix (see our response to R1 Q8), because it derived from a superseded sample and constructs not defined in the current pipeline; no claim in the revised manuscript relies on it. Analysis code and step-by-step instructions to reproduce every table and figure are provided in the replication package described in the Data availability statement.

\subsection*{Editor requirements}

\begin{itemize}
\item \textbf{Point-by-point response:} attached (this document), covering every Reviewer~1 sub-point (Q1--Q9), every Reviewer~2 point (1--9), and each editor requirement.
\item \textbf{Editable source files:} the revised manuscript is uploaded in editable source form alongside the PDF, per the production requirement.
\item \textbf{Highlights:} a separate \texttt{Highlights} file accompanies the submission with four bullet points, each verified at or below 85 characters including spaces, worded consistently with the title and the revised abstract.
\item \textbf{Graphical Abstract:} we respectfully decline the optional Graphical Abstract.
\end{itemize}

For each point below we restate the reviewer's comment (in the left column) and describe our response and the corresponding changes to the manuscript (right column), with explicit references to the revised sections and tables.

%% ------------------------------------------------------------------
\section*{Response to Reviewer 1}

\begin{longtable}{@{}L R@{}}
\toprule
\textbf{Reviewer comment} & \textbf{Response and changes to the manuscript}\\
\midrule
\endfirsthead
\multicolumn{2}{@{}l}{\footnotesize\itshape Reviewer 1 (continued)}\\
\toprule
\textbf{Reviewer comment} & \textbf{Response and changes to the manuscript}\\
\midrule
\endhead
\midrule
\multicolumn{2}{r@{}}{\footnotesize\itshape continued on next page}\\
\endfoot
\bottomrule
\endlastfoot

\textit{(Q1) ``The authors should separate the empirical question from the proposed psychological mechanism. The data can show whether users answer more after receiving help, but they do not directly measure gratitude, moral obligation, or incentive displacement.''}

\textit{(Q5) ``The `re-engagement window' interpretation is also plausible, but the paper does not observe sessions, exits, returns, or notification behavior. This should be presented as a possible explanation rather than as an established mechanism.''}

\textit{(Q7) ``The authors should acknowledge that they do not directly observe psychological mechanisms such as gratitude, obligation, or status motivation.''} &
We reframed the manuscript so the estimand is stated plainly while every psychological and session mechanism is marked as a \emph{proposed}, unmeasured interpretation; no models were re-fit for this change. The abstract states that ``the data identify short-run behavioral associations, not the psychological mechanisms (gratitude, moral obligation, incentive displacement) that theory proposes,'' and the Introduction closes that we ``test only the behavioral consequence those mechanisms predict.'' In the Theory section (``Generalized Reciprocity''; ``Reciprocity and the User Trajectory'') gratitude, moral obligation, and norm displacement are relabelled as proposed mechanisms the design cannot adjudicate. The retitled discussion subsection ``The Response-Time Pattern and Its Interpretation'' and the consolidated Limitations paragraph (Section~\ref{sec:limitations}) present the session-dynamics account as one conjecture among several. \\
\midrule

\textit{(Q1) ``The objectives are generally clear, but the rationale could be stated more modestly. The paper should clarify that its main contribution is not to establish a general theory of generalized reciprocity, but to test a short-run behavioral response to receiving an answer on Stack Overflow.''}

\textit{(Q6) ``The tenure result is potentially the most interesting contribution and should receive more emphasis than the broader claims about moral reciprocity.''}

\textit{(Q5) ``The conclusion should emphasize the narrower finding: receiving an answer is associated with a small short-run increase in helping, mainly among newer users.''} &
We revised the abstract, introduction, and conclusion without changing any estimate. The Introduction now states in its opening paragraph that ``our contribution is a better-identified short-run behavioral test on one platform,'' and positions the estimates directly against prior empirical reports of online reciprocity (Baker~2014, Chen~2019 vs.\ Wasko~2005, Joyce~2006) in the ``Empirical Challenges'' subsection. The three-contribution paragraph makes the tenure gradient the headline (``the most robust and informative pattern in the data''), and the Conclusion (Section~\ref{sec:conclusion}) leads with exactly the narrower finding the reviewer prescribes: a small pooled effect (arrival-increment HR 1.05, 95\% CI [1.04, 1.05]), concentrated among newcomers (arrival HR 1.07) and undetectable among veterans. The title's contrast (``newcomers, not veterans'') reflects this tenure-focused headline. \\
\midrule

\textit{(Q2) ``The authors should clearly distinguish questions, matched observations, interval rows, users, and help events. The reported sample sizes and event counts in the main tables are difficult to reconcile.''}

\textit{``The paper should explain more clearly how controls are duplicated or reused in matching, and how repeated observations from the same user are handled.''} &
We rewrote the ``Data'' subsection as an explicit filtering-and-matching chain: from the full non-deleted question set we removed 1{,}080 pre-launch and 8{,}029 post-cutoff cases and 2{,}671{,}121 self-answered questions, yielding 21{,}033{,}429 pre-matching questions from 4{,}764{,}048 users; after matching, the analytic file comprises 24{,}349{,}760 question-level rows (12{,}174{,}880 matched pairs) from 3{,}732{,}706 users, 50\% treated by construction, with 3{,}435{,}973 pooled helping events. The ``Propensity Score Matching'' subsection explains that matching is with replacement, so some control questions match more than one treated question --- hence more rows than unique pairs. Every table was regenerated from one pipeline run, so counts agree across the Data prose, Table~\ref{tab:desc_stats}, Table~\ref{tab:pooled_experienced} (N~$=$~24{,}349{,}760; Events~$=$~3{,}435{,}973), and Table~\ref{tab:main_results} (per-bucket N and events summing to the pooled totals). The appendix ``Implementation Details'' documents that pooled models fit a $\sim$8-million-row matched subsample (pair integrity preserved) when the interval frame exceeds the fit cap, while tenure-bucket models use all intervals. All code and reproduction instructions are in the replication package (Data availability statement). \\
\midrule

\textit{(Q3) ``Standard errors are a concern. The appendix states that robust or clustered standard errors are not used. Given repeated observations by users, repeated controls, matched pairs, and recurrent helping events, the current inference may be too optimistic.''}

\textit{``The authors should report robust standard errors clustered at least by user or matched pair, or use a bootstrap that preserves the matched structure.''} &
The installed estimator cannot deliver native clustered variances (\texttt{lifelines}' \texttt{CoxTimeVaryingFitter} raises \texttt{NotImplementedError} for \texttt{robust=True} and exposes no cluster column for time-varying fits), so we implemented the reviewer's suggested fallback: a matched-pair bootstrap that resamples whole matched pairs with replacement, refits Model~A, and takes percentile confidence intervals (100 replicates). The bootstrapped statistic is the headline treatment effect itself --- the difference-in-differences increment at answer arrival ($\beta_4$) --- so the resampling uncertainty attaches directly to the quantity the paper reports. Table~\ref{tab:pair_bootstrap} reports HR 1.05, bootstrap 95\% CI [1.03, 1.07], alongside the model-based CI [1.04, 1.05]; the bootstrap interval is also shown in the pooled regression table (Table~\ref{tab:pooled_experienced}). Relatedly, we now report the waiting-period term ($\beta_2$) as an explicit parallel-trends diagnostic with its own standard error in every table, rather than folding it into any combined contrast: because $\beta_2$ carries no sign guarantee, a summed quantity such as $\exp(\beta_2+\beta_4)$ bounds nothing, and we have removed it from the paper. We deleted the assertion that clustering was unnecessary and rewrote the ``Implementation Details'' paragraph to state that the Cox-table standard errors are model-based (inverse-Hessian), \emph{not} cluster-robust, and can be too narrow given matched pairs, reused controls, and repeated intervals, cross-referencing the bootstrap table. \\
\midrule

\textit{(Q3) ``Matching on activity history and tag-level answer rates helps, but the authors should add stronger robustness checks, such as controls for posting hour/day, question length, number of tags, user reputation, and text-based measures of question quality or complexity.''}

\textit{(Q7) ``The paper should more clearly acknowledge that answer receipt is not random and may reflect unobserved question quality or user characteristics.''} &
We re-extracted fields dropped from the processed pipeline and added the requested observable pre-treatment covariates to the propensity model --- posting hour-of-day and day-of-week, question body and title length (characters), and owner reputation --- and now match pairs \emph{exactly} on the number of tags in addition to the primary tag. Table~\ref{tab:covariates} lists these in a new ``Observable question and author characteristics'' block, and Table~\ref{tab:balance} shows they are well balanced after matching (largest residual $|$SMD$|$ 0.041, owner reputation; all below 0.05; the largest pre-matching imbalance is the tag-level accepted-answer rate, 0.695). In the pooled sample (Table~\ref{tab:observable_controls}) and the newcomer stratum (Table~\ref{tab:newcomer_robustness}) the pooled arrival increment is essentially unchanged when these controls are additionally entered into the Cox model (1.04 with controls versus 1.05 baseline), and the newcomer increment is, if anything, strengthened (1.13 versus 1.07 baseline). Limitations (Section~\ref{sec:limitations}) now states that answer receipt is non-random and that the design ``cannot cleanly separate a reciprocal response to receiving help from selection into being answered.'' \\
\midrule

\textit{(Q3) ``The response-time analysis should be modeled more flexibly. The results appear closer to an inverted-U pattern than to a simple `faster is stronger' effect.''}

\textit{(Q5) ``H3 should be rewritten or toned down. The strongest effect occurs in the 30--60 minute bin, not among the fastest responses. This does not cleanly support the original hypothesis.''}

\textit{(Q4) ``Figure 4 is useful, but the text should treat it as evidence of a non-linear response-time pattern rather than as simple support for faster responses.''} &
We replaced the single continuous log-response-time interaction with a flexible, non-parametric specification: ten discrete response-time bins, each fit as a separate Model~A (Table~\ref{tab:response_time_bins}, Figure~\ref{fig:interaction_effect}). The pattern is non-linear --- the arrival increment is near-null under 30 minutes, peaks sharply at 30--60 minutes (HR 1.19, 95\% CI [1.17, 1.21]), and is modest thereafter, with a smaller secondary elevation at multi-day delays that we report but do not build on --- an early ``sweet spot,'' not a monotone gradient. We rewrote the ``Response Time Moderation (H3)'' subsection around this specification, presenting the non-linear result first, and retitled the discussion subsection, offering H3 as \emph{qualified} support for an immediacy account rather than confirmation. The peak survives even strongly post-treatment answer-quality controls as the significant early-window maximum (1.03 [1.02, 1.04], $p<.001$; Table~\ref{tab:response_time_bins_quality}; see R2 below). \\
\midrule

\textit{(Q4) ``Table 5 should report hazard ratios and confidence intervals for both specifications, not only for the main model.''} &
We traced the missing interval to the table generator, not the model: the CI bounds for the speed specification's treatment term were computed but dropped before the table was written. We added the read-out and regenerated the table. Table~\ref{tab:pooled_experienced} now reports the coefficient, standard error, and HR with 95\% CI for the arrival increment ($\beta_4$) in both columns, together with the waiting-period diagnostic ($\beta_2$) and its standard error. The only remaining em-dash is the Main+Speed bootstrap-CI cell, because the matched-pair bootstrap (Table~\ref{tab:pair_bootstrap}) was run for Model~A alone. \\
\midrule

\textit{(Q4; see also Reviewer 2, P7) ``The authors should translate hazard ratios into absolute effect sizes, such as additional answers per 1,000 treated questions.''} &
We added an absolute-effect derivation translating each tenure-bucket hazard ratio into an absolute risk difference (ARD), number-needed-to-treat (NNT), and CI-derived NNT bounds using the observed baseline event rate (Table~\ref{tab:absolute_effects}). The newcomer arrival increment (HR 1.07) corresponds to an ARD of about 0.65 percentage points (roughly 6--7 additional helping events per 1{,}000 answered newcomer questions), NNT 154.7 (95\% CI 127.7--195.3). The absolute increment is of similar size across the active tenure range (pooled NNT $\approx$ 145, because more-tenured users' higher baseline helping rates offset their smaller relative increments) and null among the most experienced users, so the revised ``Implications for Platform Design'' rests the case for targeting newcomers on where the reciprocity response exists at all rather than on a larger absolute yield. We rewrote ``Implications for Platform Design'' to foreground this modest magnitude, restricting the practical leverage to a targeted, modest newcomer onboarding channel rather than a platform-wide multiplier. \\
\midrule

\textit{(Q7) ``The results are from Stack Overflow, a highly specific technical Q\&A community, and may not generalize to other online communities.''} &
Addressed jointly with Reviewer 2's external-validity points below. The external-validity paragraph in Limitations (Section~\ref{sec:limitations}) now leads with Stack Overflow's distinctiveness (large, relatively anonymous, gamified, technical) and refrains from claiming generalized reciprocity ``likely operates'' elsewhere, naming enterprise Slack channels, academic peer review, and intimate hobby forums as contrasting settings. \\
\midrule

\textit{(Q8) ``The theory section should be shortened and more directly connected to the empirical design.''}

\textit{``Some implementation details currently in the appendix should be moved or summarized in the main text because they affect interpretation of the results.''} &
We shortened the Theory section by roughly one-sixth (from about 1{,}550 to about 1{,}290 words), condensing the two-mechanism exposition of generalized reciprocity into a single paragraph (``Generalized Reciprocity'' subsection) and compressing the empirical-challenges discussion. The theory now links forward to the design in two places. First, the empirical-challenges subsection closes: \emph{``This confound is the central identification challenge that motivates the matched difference-in-differences design we develop in Section~\ref{sec:methods}.''} Second, a new paragraph immediately after Hypothesis~2 (end of the ``Reciprocity and the User Trajectory'' subsection) carries the design stakes forward: \emph{``If the reciprocity association is concentrated among newcomers, generalized reciprocity would function less as a general engine of cooperation among established members than as a contributor-recruitment channel operating at the point of entry, before community-specific incentives take hold --- a design implication we develop in Section~\ref{sec:design_implications}.''} We also promoted the tenure-stratified adoption figure from the appendix into the Results section (now Figure~\ref{fig:help_rate_adoption_by_tenure}, immediately after the pooled Figure~\ref{fig:help_rate_adoption_pooled}), together with its interpretive paragraph (``Response Time Moderation (H3)'' subsection), so the disaggregated evidence bearing on the session-dynamics interpretation appears in the main text. The sample-construction and event-counting sub-point is addressed under Q2/Q8 above. One transparency note: the original appendix's lifetime-contributions table was removed rather than promoted, because it derived from a superseded sample with an inconsistent data cutoff and constructs not defined in the current pipeline; no claim in the revised manuscript relies on it. \\
\midrule

\textit{(Q9) Would you recommend language editing? --- ``Yes.''} &
A full language and consistency edit was applied to the revised manuscript after the substantive revisions were complete. Representative corrections, with locations: ``before and after an user receives an answer'' now reads ``\ldots a user\ldots'' (opening of Section~\ref{sec:methods}); a garbled sentence in the setting description now reads \emph{``Its user base spans a wide range of experience levels, from learners to professional software developers''} (Section~\ref{sec:empirical_setting}); a stray parenthesis and misplaced punctuation around citations were repaired (Section~\ref{sec:empirical_setting}; ``Analytical Strategy'' subsection). Terminology and formatting were standardized throughout (``difference-in-differences,'' defined as DiD at first use in the Introduction; ``propensity score matching''; ``response-time'' as a modifier; ``hazard ratio'' defined at first use in both the abstract and the Introduction). The abstract was made self-contained by removing its cross-reference to an appendix table, and the Highlights were aligned with the revised abstract. \\

\end{longtable}

%% ------------------------------------------------------------------
\section*{Response to Reviewer 2}

\begin{longtable}{@{}L R@{}}
\toprule
\textbf{Reviewer comment} & \textbf{Response and changes to the manuscript}\\
\midrule
\endfirsthead
\multicolumn{2}{@{}l}{\footnotesize\itshape Reviewer 2 (continued)}\\
\toprule
\textbf{Reviewer comment} & \textbf{Response and changes to the manuscript}\\
\midrule
\endhead
\midrule
\multicolumn{2}{r@{}}{\footnotesize\itshape continued on next page}\\
\endfoot
\bottomrule
\endlastfoot

\textit{(Overall) ``The problem is that the topic is out-dated in the era of GAI. [\ldots] This makes the topic not interesting.''}

\textit{(P9) ``The paper treats its conclusions as timeless, but they are, at best, a pre-generative-AI baseline. The authors must add a dedicated limitations paragraph explicitly stating that their findings are a historical snapshot, not a prescriptive guide for contemporary or future platform design.''}

\textit{(P2) ``The authors overgeneralise findings to all knowledge-sharing platforms, ignoring that Stack Overflow is uniquely technical, anonymous, and status-driven and critically, its ecosystem has been structurally disrupted by generative AI since 2023.''} &
We revised ``Empirical Setting: Stack Overflow'' (Section~\ref{sec:empirical_setting}) to remove the broad-generalization claim and to state that the platform's active question-and-answer volume ``has declined sharply since the release of generative AI tools in late 2022,'' that the sample spans inception through 1 April 2025 (a period that includes but predates the steepest LLM-driven decline), and that the findings should be read as ``a historical snapshot of one large, technical, relatively anonymous, status-driven platform.'' A dedicated GAI-era paragraph in Limitations states the window is a historical snapshot, not a prescriptive guide, and cross-references the pre-2020 robustness (Table~\ref{tab:cohort_robustness}): both the pooled and newcomer associations are essentially unchanged on the pre-2020 subsample (pooled arrival HR 1.04 pre-2020 vs.\ 1.05 full sample; newcomer 1.07 in both), so the association predates rather than depends on the generative-AI era. \\
\midrule

\textit{(P3) ``The analytical sample only includes users who remained active long enough to post a question and be observed. Users who posted a question, received no answer, and immediately churned are excluded. Since this exclusion is likely correlated with the treatment (no answer) and the outcome (no future helping), the estimated treatment effect [\ldots] is probably downward-biased. The authors should conduct a sensitivity analysis using Heckman-type selection correction or, at minimum, explicitly bound the possible bias.''} &
As a point of clarification about the sampling frame: the pipeline imposes no minimum-activity, minimum-tenure, or re-observation filter, so asked-but-never-answered users are \emph{retained} as the control pool, and because each matched pair shares a fixed, pair-equalized observation window, an immediate churner simply contributes zero help events over a fully observed window. We did not implement a Heckman correction (no matching covariate provides a credible exclusion restriction, and the recurrent-event time-varying Cox is not natively Heckman-compatible); instead we added a selection-sensitivity analysis reporting transparent bounds under adversarial assumptions about zero-event controls (Table~\ref{tab:selection_bounds}): the pooled effect stays modest even under optimistic bounds, while the newcomer contrast is more churn-sensitive. The first, most-emphasized Limitations paragraph (Section~\ref{sec:limitations}) reframes the threat as right-truncation/churn and flags the newcomer sensitivity explicitly. \\
\midrule

\textit{(P4) ``The outcome variable counts only `posting an answer' as a helping event. On Stack Overflow, users frequently reciprocate through upvoting, editing, commenting, or accepting an answer. [These] behaviours require far less effort and are especially common among novices who lack expertise. By excluding these actions, the study systematically undercounts reciprocal behaviour and misclassifies many reciprocating users as non-responders. The authors should re-estimate the model using a composite helping indicator (including upvotes and comments) or justify why only answers capture the theoretically relevant construct, with reference to prior reciprocity literature.''}

\textit{(P1) ``The operationalisation of `help' is too narrow (answers only), excluding upvotes, edits, and comments --- forms of low-cost reciprocity that newcomers often use.''} &
We did both of the things the reviewer proposes. First, the justification: a new ``Why answers anchor the construct'' paragraph (Variables subsection) grounds the answers-anchored outcome in the experimental reciprocity literature --- generalized reciprocity is operationalized there as a \emph{costly} prosocial contribution (Greiner \& Levati 2005; Stanca 2009; Mujcic \& Leibbrandt 2018), and it is the willingness to bear that cost, not low-cost approval, that distinguishes reciprocity from signaling (Molm et al.\ 2007). Second, the re-estimation: we parsed the additional event types from the raw dump and re-estimated the pooled model under every broader outcome definition the data support, reported as a full decomposition in Table~\ref{tab:outcome_decomposition}: answers-only (headline) arrival HR 1.05 [1.04, 1.05]; comments-only 1.01 [1.01, 1.02]; edits-only 1.16 [1.16, 1.17]; answers$+$comments 1.02 [1.02, 1.03]; answers$+$comments$+$edits 1.05 [1.04, 1.05]. The decomposition shows the answers-only headline does not manufacture the effect: comments carry only a slight elevation, well below the answers increment, and adding them \emph{dilutes} rather than amplifies the estimate --- the opposite of the mechanical newcomer-undercount the reviewer is concerned about. Two of the channels the reviewer names cannot enter the model, for distinct reasons. Accepting an answer is a \emph{self-directed} act: the asker marks an answer on their own question, and only treated users ever receive an answer to accept, so a matched control has no counterpart behavior and the difference-in-differences contrast is degenerate rather than informative --- we note this where the decomposition is introduced. Up-votes are simply unobservable: Stack Exchange suppresses voter identity in the public dump (the voter \texttt{UserId} is null for essentially all up-vote records), so an ``up-votes given'' signal cannot be attributed to an actor on any Stack Overflow data --- a data-source limitation, not a pipeline gap. Limitations (Section~\ref{sec:limitations}) states explicitly that the reported hazard ratios describe costly contribution specifically, not a complete accounting of reciprocity. \\
\midrule

\textit{(P5) ``The paper interprets the peak reciprocity at 30--60 minutes as a `re-engagement window' effect. However, a simpler alternative is answer quality: extremely fast answers ($<$15 min) are often terse, link-only, or low-scoring, while 30--60-minute answers are typically more detailed, with code snippets and explanations, generating greater gratitude. The authors do not control for answer length, vote score, or whether the answer was accepted. Without testing this confound (e.g., by conditioning on answer quality metrics), the structural `session-dynamics' interpretation remains speculative and should be toned down.''} &
We re-extracted three observable quality signals (first-answer length in characters, vote score, and acceptance) and re-estimated the pooled speed specification, adding them one at a time as a graded decomposition (Table~\ref{tab:answer_quality_robustness}). Controlling for \emph{answer length} alone --- the reviewer's terse-versus-detailed variable, and the signal least entangled with the asker's own behavior --- barely moves the estimate (arrival 1.038 to 1.032); vote score alone brings it to 1.010; and the full specification, which additionally conditions on acceptance, pushes it below one (0.982), with the same pattern in the newcomer stratum (quality-controlled arrival 0.95; Table~\ref{tab:newcomer_robustness}). Because vote score and, especially, acceptance are realized only for answered questions and are endogenous to the outcome (accepting requires the asker to return to the site), we frame the full-control estimate as a post-treatment bad-control bound in ``Answer quality as an alternative explanation.'' and recast the session-dynamics reading as one conjecture among several (quality, session dynamics, selection). We also verified that the below-one result is not an artifact of how the vote score is coded: re-estimating the score-only and full specifications under a signed-log ($\text{asinh}$) coding and ordinal-bin codings of the score (at full data, pooled and in the newcomer stratum) leaves the full-specification increment below one under all three codings, with the better-specified signed-log and binned codings if anything pulling it marginally \emph{further} below one; the only estimate sensitive to the coding is the intermediate score-only step, which is above one under the linear coding but statistically indistinguishable from one under the alternatives (new Appendix Table~\ref{tab:score_coding_sensitivity}). We use $\text{asinh}(\text{score})$ rather than $\log(1+\text{score})$ because roughly two million answers carry net-negative (downvoted) scores, for which $\log(1+\text{score})$ is undefined. \\
\midrule

\textit{(P5, continued) ``The striking inverted-U response-time effect (peak at 30--60 minutes) conflates answer speed with answer quality; faster answers may be terse or incorrect, while 30-minute answers tend to be more thorough, a confound not tested.''} &
Beyond the pooled test above, we re-estimated each of the ten discrete response-time bins with the same three observable quality controls, fitting Model~A plus those signals per bin so the baseline and quality columns are directly comparable (Table~\ref{tab:response_time_bins_quality}). The 30--60-minute peak attenuates from 1.19 [1.17, 1.21] to 1.03 [1.02, 1.04] but remains the clearly significant ($p<.001$) early-window maximum, while the same controls push the fastest bins significantly \emph{below} one (0.95) --- a signature of over-adjustment, since score and acceptance are strongly post-treatment (acceptance requires the asker's return, part of the re-engagement pathway that produces the outcome). That the peak survives an adjustment biased against it indicates that answer timing retains explanatory content beyond observable quality. ``Response Time Moderation (H3)'' reports this as genuine but bounded support for the immediacy account; a persistent secondary elevation at one to three days (1.09) is left uninterpreted because the controls are post-treatment. \\
\midrule

\textit{(P6) ``The platform has undergone massive cultural, technical, and policy shifts from 2008 to 2025 --- including reputation inflation, spam filters, tag evolution, and most critically, the rise of ChatGPT post-2022. [\ldots] the tenure-stratified analysis pools newcomers from 2008 with those from 2024 as if they share the same `newcomer' experience. This ignores cohort heterogeneity. The authors should include year-of-entry fixed effects or conduct robustness checks on pre-2020 subsamples to isolate the pure human-only reciprocity mechanism.''} &
A year-of-entry fixed-effects specification would be collinear with our matched design --- calendar year already enters as an exact-matching stratum, so registration cohort is absorbed once question-year is matched --- so we instead pursued the pre-2020 subsample the reviewer also proposed, both pooled and in the newcomer stratum. Re-estimating Model~A before 2020 (Table~\ref{tab:cohort_robustness}) leaves the arrival increment essentially unchanged (pre-2020 pooled arrival HR 1.04 versus 1.05 full sample; pre-2020 newcomer 1.07 versus 1.07 full sample). We report this in the ``Experience Moderation of Reciprocity (H2)'' subsection and frame the findings as a pre-/early-generative-AI snapshot. \\
\midrule

\textit{(P7) ``The paper does not report the Number Needed to Treat (NNT) --- i.e., how many newcomers must receive an answer to generate one additional helping event. Without this, platform designers cannot assess whether investing in faster response systems yields meaningful returns. The authors should calculate and report NNT, and revise their `Implications for Design' section to reflect modest effect sizes.''} &
Addressed jointly with Reviewer 1's absolute-effects request above. Table~\ref{tab:absolute_effects} reports ARD, NNT, and NNT confidence bounds by tenure bucket (newcomer NNT 154.7, 95\% CI 127.7--195.3; pooled NNT $\approx$ 145; similar absolute yields, with the design case resting on the null response among veterans), and ``Implications for Platform Design'' foregrounds the modest absolute magnitude. \\
\midrule

\textit{(P8) ``Posting a question increases a user's profile visibility: other members click on their profile, view their history, and may trigger a sense of social obligation or reputation concern. This `exposure effect' is not equivalent to gratitude for receiving an answer, yet it operates precisely during the post-answer period and is not separately identified. The waiting-period covariate absorbs baseline question-asking effects, but it does not capture the increased visibility that persists after an answer arrives. A placebo test using questions that received no answers but high views could help rule out this confound; currently, it is not attempted.''} &
We added a dedicated exposure/profile-visibility paragraph to Limitations (Section~\ref{sec:limitations}) that names the confound, states the waiting-period term (\textit{treatedPostQuestion}) absorbs the act of asking but not residual post-answer visibility, and frames the issue as a threat to mechanism \emph{attribution} rather than to the existence of the behavioral pattern. We explain why the suggested view-count placebo is not identified with the archival data available to us: the dump records \textit{ViewCount} only as a single cumulative snapshot at the dump date, conflating a question's age, topic popularity, and years of later search-engine traffic with the contemporaneous exposure the placebo would require. We softened the attribution language in the ``Identification.'' paragraph and in the H1 results (the small, negative pooled waiting-period coefficient ($\beta_2 = -0.007$, opposite in sign to the arrival increment) is stated to rule out only an upward drift from the act of asking, with a forward pointer to Limitations). \\
\midrule

\textit{(Q9) Would you recommend language editing? --- ``Yes.''} &
Addressed jointly with Reviewer 1's Q9 above: a full language and consistency edit was applied to the revised manuscript after the substantive revisions were complete. \\

\end{longtable}

We hope these revisions address the reviewers' concerns. We are grateful for feedback that materially improved the identification, the transparency of the sample and inference, and the calibration of our claims, and we are happy to provide any further analyses the reviewers consider useful.

\end{document}
